\documentclass[11pt]{article}
\usepackage{fontspec}
\setmainfont{Times New Roman}
\setsansfont{Arial}
\setmonofont{Consolas}
\usepackage{amsmath,amssymb,mathtools}
\usepackage{geometry}
\usepackage{microtype}
\newcommand{\mo}{\mathfrak{o}}
\newcommand{\mO}{\mathfrak{O}}
\newcommand{\mT}{\mathfrak{T}}
\newcommand{\mP}{\mathfrak{P}}
\newcommand{\mpideal}{\mathfrak{p}}
\geometry{a4paper,margin=2.35cm}
\setlength{\parindent}{1.25em}
\setlength{\parskip}{0pt}
\makeatletter
\renewcommand\footnotesize{\@setfontsize\footnotesize{7.4}{8.4}\selectfont}
\makeatother
\emergencystretch=100em
\allowdisplaybreaks
\sloppy
\hbadness=10000
\vbadness=10000
\hfuzz=2pt
\vfuzz=10pt
\raggedbottom

\begin{document}

% Unit: Paper31 Section08 Entry09 general discriminant theorem statement v001. Direct Interslavic Latin translation from clean RA34 line 445 / segment P31-S0102, checked against printed scan page 103 / PDF page 22 and Claude handoff OCR line 433. No new footnote belongs to this theorem-statement unit.
\setcounter{footnote}{39}

\noindent 5. \textbf{Obči teorem o diskriminantu.} \emph{Prosty ideal $\mpideal$ iz multiplikacijskogo kolca $\mo$ vhodi v diskriminantny ideal porjadka vzhodno k $\mo$ togda i samo togda, kogda v razloženju ideala $\mT\mpideal$ na po parah vzajemno proste primarne komponenty iz $\mT$ pojavja se najmanje jedna vlastno primarna komponenta abo najmanje jeden prosty ideal drugogo roda.}

\end{document}
